\clearpage
\section{Systemkomponenten}

Innerhalb der Systemgrenze werden drei Komponenten unterschieden. Diese Aufteilung folgt den
drei fachlichen Aufgaben: Kommunikation ausführen, Kundenantwort erfassen und Ergebnisse für
den Disponenten aufbereiten.

\begin{figure}[H]
\centering
\resizebox{0.98\textwidth}{!}{%
\begin{tikzpicture}[
  font=\small,
  component/.style={draw=MinovaBlue,thick,rounded corners=3pt,minimum width=3.5cm,minimum height=1.25cm,align=center,fill=SoftBlue},
  external/.style={draw=MinovaDark,rounded corners=3pt,minimum width=2.7cm,minimum height=1.2cm,align=center,fill=SoftGray},
  arrow/.style={-{Latex[length=2.4mm]},thick}
]
  \node[external] (dispo) at (0,0) {\textbf{DISPO}};
  \node[component] (service) at (5.0,0) {\textbf{Avisierungsservice}\\Kommunikation und Rückmeldung};
  \node[component] (page) at (10.1,1.35) {\textbf{Bestätigungsseite}\\Kundenansicht};
  \node[component] (dashboard) at (10.1,-1.35) {\textbf{Avisierungsdashboard}\\Disponentenansicht};
  \node[external] (customer) at (14.5,1.35) {\textbf{Kunde}};
  \node[external] (dispatcher) at (14.5,-1.35) {\textbf{Disponent}};

  \draw[arrow] ([yshift=0.22cm]dispo.east) -- ([yshift=0.22cm]service.west);
  \draw[arrow] ([yshift=-0.22cm]service.west) -- ([yshift=-0.22cm]dispo.east);
  \draw[arrow] (service) -- (page);
  \draw[arrow] (service) -- (dashboard);
  \draw[arrow] (page) -- (customer);
  \draw[arrow] (dashboard) -- (dispatcher);

  \begin{scope}[on background layer]
    \node[draw=MinovaBlue,very thick,rounded corners=5pt,fit=(service)(page)(dashboard),inner sep=0.38cm,label={[font=\bfseries,text=MinovaBlue]above:System zur digitalen Lieferavisierung}] {};
  \end{scope}
\end{tikzpicture}
}
\caption{Komponenten und fachliche Schnittstellen}
\label{fig:components}
\end{figure}


\subsection{Avisierungsservice}

Der Avisierungsservice ist die zentrale Komponente der Lösung. Er übernimmt die Daten des
geplanten Auftrags, erzeugt die Avisierung, veranlasst den Nachrichtenversand und speichert
die Antwort des Kunden. Au\ss{}erdem stellt er die Informationen für Bestätigungsseite und
Dashboard bereit und meldet das Ergebnis an DISPO zurück.

\subsection{Bestätigungsseite}

Die Bestätigungsseite ist die für den Kunden bestimmte Web-Oberfläche. Sie zeigt die
wesentlichen Lieferdaten und bietet zwei eindeutige Aktionen: Termin bestätigen oder Termin
ablehnen. Bei einer Ablehnung kann der Kunde einen Kommentar hinterlassen. Der Zugriff erfolgt
über den persönlichen Link aus der Benachrichtigung.

\subsection{Avisierungsdashboard}

Das Avisierungsdashboard ist die interne Arbeitsoberfläche für Disponenten. Es fasst
Kundenrückmeldungen zusammen, hebt Fälle mit Handlungsbedarf hervor und zeigt in der
Detailansicht Kommentare und den bisherigen Verlauf. Von dort kann eine unveränderte Anfrage
über einen anderen Kanal erneut versendet oder eine telefonische Bestätigung manuell erfasst
werden.
